\section{Discussion}
\label{sec:discussion}

The design separates two questions that are easy to run together: which quantity does a sentence
report, and is the number right. The first needs reading; the second needs arithmetic. In \pv{} a
language model answers only the first and code answers the second. On this manuscript that
division gave no false alarms with a strong model, and with weaker models it turned wrong answers
into lost coverage rather than wrong conclusions.

The division does not remove every risk. When a wrong cell holds the same rounded value as the
right one, the value check passes. The paper renders correctly, but after a rerun the text would
follow the wrong cell. The DockPrepBench tables contain many such coincidences, because rates
round to one decimal and the same rate recurs across engines and arms. Two safeguards are
available: review migrations before accepting them, or accept only addresses on which two models
agree, which on this manuscript left one disagreement with Claude in 110 numbers. The benchmark
shows why one of them is needed: the most precise local model in shortlist mode still chose a
wrong cell with the right value for 15 numbers. The benchmark recorded only totals per model, so
whether agreement between two local models removes those errors is not yet measured.

The evaluation has clear limits. It uses one manuscript, written by the author of the tool, in one
field, with tables designed for it. The corruptions are synthetic and of one kind, a change in the
last digit; a number attached to the wrong arm, or a sign error, could behave differently. Integers
were left out. The shortlist comparison uses Claude's answers as the reference before an
independent gold set existed, and the gold set that now exists comes from the same migration.
Whether these results hold on other papers is untested. A corpus of manuscripts by other authors,
with their analysis output, is the obvious next step, and is needed before the approach can be
recommended for general use.

\pv{} also asks something of the author: the analysis has to write its results as tables with
stable keys. Many projects already do this, but a number computed in a notebook cell and pasted
into the text cannot be traced until it is written to a table. In our experience this is a useful
constraint, since the 28 untraceable numbers found by the audit were exactly the values most
likely to go stale.
